\chapter{需求分析}

\section{需求获取方法}

本系统的需求通过五种方法获取，覆盖从用户行为分析到工程质量审查的完整链路。

\textbf{真实账单处理流程分析}：项目启动阶段，团队收集支付宝与微信支付的原始账单文件，模拟用户从下载账单、整理数据到生成报表的完整操作流程，提炼出导入、解析、分类、分析、报表五大核心环节，作为功能需求的主要依据。

\textbf{小组讨论与范围界定}：团队依据开发周期与技术能力，对需求进行优先级排序，将核心闭环（认证、导入、分类、分析、报表）确定为最高优先级，将消费人格分析、家庭组织、模型切换等纳入扩展范围，形成清晰的 MVP 边界。

\textbf{误分类案例分析}：团队收集了多种语义误判场景——如“东南大学饭卡充值”被通用分类归为“学习”、连锁便利店消费被归为“购物”而非“餐饮”——驱动了大模型初分、分类理由展示和人工校正工作台的需求。系统需在给出分类建议的同时，提供置信度与推理理由，允许用户在校正工作台中批量修正。

\textbf{原型与联调反馈}：在前后端迭代过程中，前端页面的交互细节和后端接口的返回结构经过多轮联调调整。例如，导入页面的进度轮询机制、校正工作台的左右分栏布局、家庭组织的角色权限控制，均来自原型验证阶段的反馈。

\textbf{工程质量审查}：团队对照安全、隐私、性能、可部署性、可维护性、可测试性等质量维度，识别并补齐非功能需求。这一步骤确保系统不仅是功能可用的原型，也是具备工程质量的交付物。

\headingneedspace{19}
\section{利益相关者分析}

系统涉及四类利益相关者，其关注点与需求来源如表所示。

\begin{table}[H]
\centering
\caption{利益相关者分析}
\label{tab:stakeholders}
\begin{tabular}{p{3.2cm} p{5.2cm} p{5.2cm}}
\toprule
\textbf{利益相关者} & \textbf{关注点} & \textbf{需求来源} \\
\midrule
普通用户 & 上传账单、查看 AI 分类结果、修正语义误判、查看消费画像、生成报表 & 个人记账与财务分析场景 \\
\hline
家庭组织管理员 & 创建家庭组织、添加成员、管理角色、查看家庭聚合分析 & 家庭账本协作场景 \\
\hline
系统运维人员 & 配置模型供应商、管理重试队列、部署与监控系统 & 系统部署与运维场景 \\
\hline
开发与维护人员 & 实现系统功能、控制开发范围、维护代码质量、准备后续扩展 & 工程开发与维护过程 \\
\bottomrule
\end{tabular}
\end{table}

\section{功能需求}

系统共识别 15 项功能需求，按优先级分为 P0（核心闭环必备）、P1（完整性与体验增强）和 P2（未来改进）。当前实现覆盖全部 P0 与 P1 需求。

\begingroup
\small
\begin{longtable}{p{1.0cm} p{3.5cm} p{0.8cm} p{6.2cm}}
\caption{功能需求列表}
\label{tab:functional-requirements}\\
\toprule
\textbf{ID} & \textbf{需求描述} & \textbf{优先级} & \textbf{实现依据} \\
\midrule
\endfirsthead
\multicolumn{4}{c}{\small 表~\thetable\quad 功能需求列表（续）} \\
\toprule
\textbf{ID} & \textbf{需求描述} & \textbf{优先级} & \textbf{实现依据} \\
\midrule
\endhead
\midrule
\multicolumn{4}{r}{\small 接下页} \\
\endfoot
\bottomrule
\endlastfoot
FR-01 & 用户注册、登录及 Token 鉴权 & P0 & \texttt{routes\_auth.py}、前端登录页、JWT 中间件 \\
\hline
FR-02 & 多用户数据隔离，账单与报表按用户归属 & P0 & \texttt{dependencies.py} 依赖注入、资源归属校验 \\
\hline
FR-03 & 支持微信/支付宝账单文件上传导入 & P0 & \texttt{routes\_imports.py}、导入页面 \\
\hline
FR-04 & 不同平台账单统一解析为标准化交易结构 & P0 & \texttt{services/parsers.py}、\texttt{test\_parsers.py} \\
\hline
FR-05 & 大模型语义分类，记录来源、置信度与理由 & P0 & \texttt{services/classifiers.py}、\texttt{test\_classifiers.py} \\
\hline
FR-06 & 交易列表筛选（日期、平台、分类、关键词）与分页 & P0 & \texttt{routes\_transactions.py}、交易列表页面 \\
\hline
FR-07 & 人工校正低置信度或语义误判交易 & P0 & \texttt{routes\_transactions.py} PATCH 接口、校正工作台 \\
\hline
FR-08 & 用户自定义分类管理 & P1 & \texttt{routes\_categories.py}、分类管理页面 \\
\hline
FR-09 & Dashboard 展示收入、支出、结余、储蓄率、分类占比 & P0 & \texttt{routes\_dashboard.py}、\texttt{services/analytics.py} \\
\hline
FR-10 & 消费人格画像、财务健康评分与问卷自评 & P1 & \texttt{routes\_personality.py}、\texttt{services/personality.py} \\
\hline
FR-11 & 家庭组织创建与成员角色管理 & P1 & \texttt{routes\_organizations.py}、\texttt{services/organizations.py} \\
\hline
FR-12 & 家庭组织 Dashboard 与消费人格聚合 & P1 & \texttt{organization\_id} 查询参数、组织页面 \\
\hline
FR-13 & PDF 报表生成与下载，支持个人与家庭模式 & P0 & \texttt{routes\_reports.py}、\texttt{services/reports.py} \\
\hline
FR-14 & 分类供应商切换（规则/外部 API/本地模型） & P1 & \texttt{services/classifiers.py} CompositeClassifier \\
\hline
FR-15 & 外部模型失败重试队列与管理员控制 & P1 & \texttt{routes\_classification.py}、\texttt{retry\_queue.py} \\
\end{longtable}
\endgroup

\headingneedspace{6}
FR-01 至 FR-07、FR-09、FR-13 构成系统的核心闭环：认证后导入账单并自动分类，经人工校对后产出 Dashboard 分析与 PDF 报表。FR-08、FR-10 至 FR-12、FR-14、FR-15 在此闭环上增加了消费人格画像、家庭协作和运维可控性。

\section{非功能需求}

系统定义了 8 项非功能需求，涵盖安全性、隐私保护、可用性、性能、可部署性、可维护性、可测试性与可扩展性。

\begin{table}[H]
\centering
\caption{非功能需求列表}
\label{tab:nfr}
\begin{tabular}{p{1.5cm} p{3cm} p{1cm} p{7.8cm}}
\toprule
\textbf{ID} & \textbf{需求名称} & \textbf{优先级} & \textbf{说明} \\
\midrule
NFR-01 & 安全性 & P0 & 密码采用哈希存储，业务接口均需 JWT 鉴权，资源访问按用户归属校验 \\
\hline
NFR-02 & 隐私保护 & P0 & 账单数据、分类结果与 PDF 报表属于敏感个人财务数据，用户间严格隔离 \\
\hline
NFR-03 & 可用性 & P0 & 用户可按“导入→初分→校正→分析→报表”路径完成核心任务；家庭功能不破坏个人模式 \\
\hline
NFR-04 & 性能 & P1 & 非模型调用接口（认证、查询、聚合分析）在单用户正常操作下 P95 响应时间应低于 1 秒；大模型分类调用设有 60 秒超时保护（\texttt{model\_timeout\_seconds=60}）与最多 2 次自动重试；≤1000 笔交易的单批次文件解析阶段应在 10 秒内完成，进度通过轮询接口实时反映 \\
\hline
NFR-05 & 可部署性 & P0 & 支持本地开发模式、Docker Compose 部署与裸机部署三种方式 \\
\hline
NFR-06 & 可维护性 & P0 & 前后端分离；后端按 routes/services/models/schemas 分层组织 \\
\hline
NFR-07 & 可测试性 & P0 & 后端具备 pytest 自动化测试套件，前端通过 TypeScript 类型检查与生产构建验证 \\
\hline
NFR-08 & 可扩展性 & P1 & 分类供应商、报表生成、异步任务与部署组件预留扩展接口 \\
\bottomrule
\end{tabular}
\end{table}

\section{用户场景分析}

\subsection{普通用户场景}

普通用户的核心场景覆盖“导入—分类—校正—分析—报表”的完整链路：

\begin{enumerate}
  \item 用户注册账号并登录系统，进入应用主界面。
  \item 进入账单导入页面，上传微信或支付宝账单文件（CSV 格式）。
  \item 系统创建导入批次，后台解析账单并将不同平台格式转换为统一交易记录。
  \item 系统调用大模型对每条交易进行语义初分，记录分类来源、置信度与分类理由。
  \item 用户在交易列表中按日期、平台、分类、关键词筛选查看交易。
  \item 在分类校正工作台中，用户查看待校正交易队列，对低置信度或语义误判的分类进行人工修正。
  \item 用户进入 Dashboard 查看收入支出概览、储蓄率、分类占比图与 Top 商户。
  \item 用户进入消费人格页面，查看基于交易数据计算的消费人格画像、财务健康评分，并可通过问卷进行自评对比。
  \item 用户进入报表中心，选择日期范围与导入文件，生成并下载 PDF 报表。
\end{enumerate}

\subsection{家庭组织场景}

家庭组织场景描述管理员创建家庭、管理成员并查看聚合分析的过程：

\begin{enumerate}
  \item 用户创建家庭组织，系统自动将创建者设为 owner 角色。
  \item owner 或 admin 通过用户名搜索添加成员，并可提升普通成员为 admin。
  \item 家庭成员进入组织页面查看家庭合计收入、支出、交易数与消费人格聚合雷达图。
  \item owner 或 admin 可触发家庭范围内的重分类操作；普通 member 只能查看聚合结果。
  \item 不传入 \texttt{organization\_id} 参数时，Dashboard、Personality 和 Reports 仍按个人账本模式工作，保持向后兼容。
\end{enumerate}

\subsection{系统运维场景}

运维人员关注系统部署、配置与故障恢复：

\begin{enumerate}
  \item 运维人员配置后端运行环境，包括数据库连接、Redis 缓存、模型供应商与 API 密钥。
  \item 部署前端静态资源、Nginx 反向代理、FastAPI 后端、PostgreSQL 与 Redis。
  \item 当模型请求超时或失败时，运维人员通过重试状态接口查看队列中的交易数量与供应商分布。
  \item 运维人员通过管理界面将历史超时或失败交易重新放入重试队列，恢复分类流程。
\end{enumerate}

\section{异常场景处理}

系统对以下异常场景定义了明确的预期行为：

\begin{table}[H]
\centering
\caption{异常场景与系统预期行为}
\label{tab:exception-scenarios}
\begin{tabular}{p{4.8cm} p{8.5cm}}
\toprule
\textbf{异常场景} & \textbf{系统预期行为} \\
\midrule
用户未登录访问业务接口 & 返回 401 未授权状态码，前端跳转至登录页 \\
\hline
用户上传不支持的账单格式 & 导入失败并记录具体错误信息，页面提示用户重新上传 \\
\hline
重复导入同一笔交易 & 通过 \texttt{dedupe\_hash} 字段去重，避免重复记录入库 \\
\hline
模型供应商不可用 & 交易进入重试队列，后台串行重试；超时耗尽后标记为 retry\_failed \\
\hline
用户尝试访问他人报表或交易 & 后端按 \texttt{current\_user.id} 校验资源归属，拒绝非授权访问 \\
\hline
报表生成失败 & 记录任务状态与错误信息，用户可调整参数后重新触发生成 \\
\bottomrule
\end{tabular}
\end{table}

\section{用例建模}

系统用例图展示三类参与者（普通用户、家庭管理员、系统运维人员）与系统功能之间的关系。普通用户可操作认证、导入、查看交易、校正分类、管理自定义分类、查看 Dashboard、生成报表、查看消费人格等用例；家庭管理员额外可创建组织和管理成员；系统运维人员负责供应商配置与重试队列管理。导入账单用例包含大模型语义分类作为扩展关系。

\begin{figure}[H]
  \centering
  \includegraphics[width=0.92\textwidth]{用例图.png}
  \caption{系统用例图}
  \label{fig:usecase}
\end{figure}

\section{用例规约}

以下对两个核心用例给出结构化用例说明，包含基本流程、替代流程与异常流程。

\subsection{UC-01：导入账单与语义分类}

\begin{table}[H]
\centering
\caption{用例规约：导入账单与语义分类}
\label{tab:uc01}
\small
\begin{tabular}{p{2.8cm} p{10.5cm}}
\toprule
\textbf{用例名称} & 导入账单与语义分类 \\
\midrule
\textbf{参与者} & 普通用户（主）；大模型分类服务（外部系统） \\
\hline
\textbf{触发条件} & 用户进入账单导入页面，选择账单文件并点击上传 \\
\hline
\textbf{前置条件} & 用户已登录；待上传文件为微信或支付宝 CSV 格式 \\
\hline
\textbf{后置条件} & 交易记录入库完毕，语义初分完成，批次状态置为 \texttt{done} 或 \texttt{partial\_failed} \\
\hline
\textbf{基本流程} &
1）用户拖拽或选择账单文件。\newline
2）前端调用 \texttt{POST /api/imports} 上传文件内容。\newline
3）系统创建 \texttt{ImportBatch}，将原始文件写入磁盘。\newline
4）后台线程解析 CSV，将各平台格式转换为统一 \texttt{Transaction} 结构。\newline
5）对每笔新交易计算 \texttt{dedupe\_hash}，跳过已存在的重复记录。\newline
6）调用 \texttt{CompositeClassifier}（缓存→外部 API→本地模型→规则回退）进行语义分类，结果写入 \texttt{auto\_category\_id}、\texttt{auto\_confidence}、\texttt{auto\_reason}。\newline
7）批次状态更新为 \texttt{done}，前端轮询后展示完成提示。\\
\hline
\textbf{替代流程 A} & 部分交易模型超时（超过 60 秒）：失败交易标记为 \texttt{retry\_queue}，其余交易正常入库，批次状态置为 \texttt{partial\_failed}。 \\
\hline
\textbf{替代流程 B} & 文件含已导入交易：\texttt{dedupe\_hash} 自动去重，跳过重复记录，批次仍正常完成。 \\
\hline
\textbf{异常流程} & 文件格式无法识别（非微信／支付宝 CSV）：文件状态标记为 \texttt{failed}，错误信息写入 \texttt{error\_message}，页面提示用户重新上传。 \\
\bottomrule
\end{tabular}
\end{table}

\subsection{UC-02：人工校正分类结果}

\begin{table}[H]
\centering
\caption{用例规约：人工校正分类结果}
\label{tab:uc02}
\small
\begin{tabular}{p{2.8cm} p{10.5cm}}
\toprule
\textbf{用例名称} & 人工校正分类结果 \\
\midrule
\textbf{参与者} & 普通用户 \\
\hline
\textbf{触发条件} & 用户进入分类校正工作台 \\
\hline
\textbf{前置条件} & 系统中存在 \texttt{needs\_review=true} 的交易记录 \\
\hline
\textbf{后置条件} & 目标交易的 \texttt{final\_category\_id} 已更新，\texttt{needs\_review} 置为 \texttt{false}，校正结果实时反映在 Dashboard 及报表数据中 \\
\hline
\textbf{基本流程} &
1）前端加载待校正队列（\texttt{needs\_review=true}，按置信度升序排列）。\newline
2）用户选中一笔交易，右侧面板展示商户、金额、当前分类、置信度与模型推理理由。\newline
3）用户从系统分类或自定义分类中选择正确类别，点击"确认校正"。\newline
4）前端调用 \texttt{PATCH /api/transactions/\{id\}}，提交新的 \texttt{final\_category\_id}。\newline
5）后端更新交易记录，将 \texttt{needs\_review} 置为 \texttt{false}。\newline
6）前端从待校正队列移除该交易，自动加载下一笔。\\
\hline
\textbf{替代流程} & 用户认为当前分类正确：直接点击"确认当前分类"，后端同样将 \texttt{needs\_review} 置为 \texttt{false}，\texttt{final\_category\_id} 保留原 \texttt{auto\_category\_id}。 \\
\hline
\textbf{异常流程} & 所需自定义分类尚不存在：用户跳转至分类管理页面新增分类后，返回工作台重新操作。 \\
\bottomrule
\end{tabular}
\end{table}

\section{业务流程建模}

核心业务活动图描述了从用户登录到报表生成的端到端流程。流程始于用户登录，接着上传账单文件，系统创建导入批次后进行解析和去重；成功解析的交易经过大模型语义初分，高置信度结果直接入库，低置信度交易进入人工复核队列；用户完成校正后，交易数据进入 Dashboard 聚合分析，随后流向消费人格与财务健康分析模块；在报表生成环节，系统根据是否选择家庭组织参数，分别生成个人 PDF 报表或家庭聚合 PDF 报表。

\begin{figure}[!htbp]
  \centering
  \includegraphics[
    width=0.58\textwidth,
    height=0.47\textheight,
    keepaspectratio
  ]{核心业务活动图.png}
  \caption{核心业务活动图}
  \label{fig:activity}
\end{figure}
